%==============================================================
% BAB V  SIMPULAN DAN SARAN
%==============================================================

\chapter{SIMPULAN DAN SARAN}


%--------------------------------------------------------------
\section{Simpulan}
%--------------------------------------------------------------

Simpulan merupakan jawaban dari tujuan yang sudah ditentukan dan tidak
dimaksudkan sebagai ringkasan hasil.
Dalam Simpulan, penulis harus dan hanya menjawab masalah dan tujuan penelitian
yang telah dirumuskan pada Pendahuluan.
Simpulan merupakan generalisasi dari hasil penelitian dan argumentasi penulis,
atau pernyataan singkat yang merupakan hakikat dari bab Hasil dan Pembahasan
atau hasil pengujian berbagai hipotesis yang berkaitan.

Simpulan merupakan hasil penelitian yang boleh jadi telah dikemukakan dalam
perumusan masalah dan telah diberi jawaban sementara berupa hipotesis.
Dalam menulis simpulan, penulis harus membedakan dugaan, temuan, dan simpulan
hasil studi.
Pernyataan simpulan harus dilakukan secara cermat dan hati-hati.
Penyampaian simpulan ini dapat dilakukan sebanyak 3 kali, yakni dalam
Pembahasan, Simpulan, dan Abstrak sehingga diperlukan kecermatan untuk
menyajikannya dengan ungkapan yang berbeda-beda.


%--------------------------------------------------------------
\section{Saran}
%--------------------------------------------------------------

Saran seyogianya mengarah ke implikasi atau tindakan lanjutan yang harus
dilakukan sehubungan dengan temuan atau simpulan penulis.
Saran yang dikemukakan harus berkaitan dengan pelaksanaan atau hasil penelitian.
Dengan demikian saran ini mengemukakan hal-hal yang perlu diteliti lebih lanjut
terutama untuk memperbaiki kelemahan atau kekurangan dalam penelitian yang
dilakukan atau perbaikan asumsi yang diambil sehingga didapatkan hasil yang
lebih baik.
Jadi, saran tersebut harus diuraikan secara spesifik.
Jangan menyarankan hal-hal yang tidak dianalisis dan dibahas dalam penelitian
serta terkesan menggurui atau memuaskan keinginan peneliti.
Untuk penelitian yang berkaitan dengan permasalahan kebijakan, tidak perlu
menyarankan kebijakan yang tidak berkaitan dengan hasil penelitian.
